\documentclass[letterpaper, 12pt]{article}
\usepackage[top=2 cm, bottom=2 cm,left=1.5 cm, right=1.5 cm]{geometry}

\usepackage[spanish,es-nodecimaldot]{babel}
\usepackage[utf8]{inputenc}
\usepackage{csquotes}

\setlength{\parindent}{1cm}
\pagestyle{empty}

\usepackage[colorlinks=true, urlcolor=blue, linkcolor=red]{hyperref} % Para insertar hipervínculos
\usepackage{amsmath}
\usepackage[outputdir=aux]{minted}
\usepackage{amssymb}
\usepackage{multicol}
\usepackage{enumerate}
\usepackage{graphicx} % Incluir imágenes
\usepackage[document]{ragged2e}
\usepackage{mathrsfs}
\usepackage[T1]{fontenc}
\usepackage[table,xcdraw]{xcolor}
\usepackage{float} %required for the placement specifier H
\makeatletter
\makeatother
% -----------------------------------------------

% Preámbulo

\title{Tarea 3: Divide y vencerás}
\author{Christopher Gómez}
\date{Enero 2024}

\begin{document}
\parbox[t]{.5\linewidth}{
    \centering
    \includegraphics[scale=0.4]{logo.png}
    \begin{center}
        UNIVERSIDAD SIMÓN BOLÍVAR \\
        CI-5651 - Diseño de Algoritmos I \\
        Prof. Ricardo Monascal \\
    \end{center}
}
\hfill \framebox[5.5cm][c]{
        \parbox[t]{.8\linewidth}{
        \centering
       Resuelto por:\\ Christopher Gómez
        }}

\phantom{This text will be invisible} \\
\centerline {\textbf{Tarea 3: Divide y vencerás}}
\justify

\begin{enumerate}

% -------- PREGUNTA 1 --------

\item Para las siguientes recurrencias, use el teorema maestro visto en clase para
hallar una fórmula cerrada para $T(n)$:

\begin{enumerate}
    \item $T(n) = 3T(\frac{n}{4}) + \frac{3(n^2 - 1)}{2}$
    \item $T(n) = 7T(\frac{n}{7}) + 2n - 3$
    \item $T(n) = 5T(\frac{n}{2}) + 2n$
    \item $T(n) = \frac{\sum_{i=1}^n (T(\frac{n}{2}) + i)}{n}$
\end{enumerate}

% -------- PREGUNTA 2 --------

\item Definimos los números de Perrin con la siguiente recurrencia:

% Funcion a trozos
\[
P(n) =
\begin{cases}
    3 & \text{si } n = 0 \\
    0 & \text{si } n = 1 \\
    2 & \text{si } n = 2 \\
    P(n - 2) + P(n - 3) & \text{si } 3 \leq n
\end{cases}
\]

Se desea que diseñe un algoritmo que tome tiempo $\Theta(\log n)$ para hallar el valor de $P(n)$. Puede suponer que todas las operaciones aritméticas involucradas (sumas, multiplicaciones, etc.) toman tiempo $\Theta(1)$.


% -------- PREGUNTA 3 --------

\item Dada una cadena de caracteres $S[1..n]$ compuesta únicamente de paréntesis que abren y que cierran, queremos un árbol de segmentos para hacer consultas de la forma $maxBP(i, j)$ que debe devolver la longitud de la sub-cadena (no necesariamente contigua) bien parentizada más larga, comprendida en el rango $[i..j]$.

Por ejemplo, si $S = [(, ), ), (, (, ), ), (, (, ), ), (]$ y solicitamos $maxBP(3, 10)$, el rango corresponde a $S' = [), (, (, ), ), (, (, )]$.

La sub-cadena bien parentizada de mayor longitud en $S'$ sería $[(, (, ), )]$ (de longitud 4).

Describa cómo sería el proceso con el cuál construiría el valor precalculado de cada nodo,
incluyendo el caso base (las hojas) y el caso recursivo (nodos intermedios)


\end{enumerate} \vspace{4mm}

\end{document}
